\documentclass[11pt]{article}
\usepackage[margin=0.9in]{geometry}
\usepackage{graphicx,amsmath,amssymb,booktabs,hyperref,xcolor}
\graphicspath{{/tmp/cheby_h0/figs/gamma_sweep/}}
\hypersetup{colorlinks=true,linkcolor=blue,urlcolor=blue}
\setlength{\parskip}{6pt}
\title{K=3 CRRA REE: $\gamma$-sweep with strict-$h{=}0$ machine-$\varepsilon$ FPs}
\author{Fixed-Point Factory --- $\gamma$-sweep at $\tau{=}1$, $G{=}7$}
\date{June 2026}

\begin{document}
\maketitle

\section*{Summary}

Sweep of the risk-aversion parameter $\gamma \in \{0.1, 0.3, 0.5, 0.8, 1.0,
1.5, 2.0, 3.0, 5.0\}$ at fixed $\tau=1$, on the $G{=}7$ Lobatto--atanh grid.
For each $\gamma$ we run the strict-$h{=}0$ POU+chebroots operator
auto-selecting the best $N_Q$ from $\{12,14,16,18,20,22\}$, and nail the
fixed point with Anderson-accelerated Picard + Newton--Krylov.

\textbf{Methodology recap.}
\begin{itemize}
  \item \textbf{Operator}: strict $h{=}0$, tensor-Chebyshev representation
        of $P(u_1,u_2,u_3)$ on Lobatto--atanh nodes (stretching $c{=}2$).
        Per-cell co-area integral via partition-of-unity over both axes,
        all real roots from chebroots (companion-matrix eigvals in numba).
        $G_p{=}121$ logit-uniform table for $\mu(p, u_k)$ lookup.
  \item \textbf{Solver}: Anderson Type-II with memory $m{=}10$, then
        SciPy \texttt{newton\_krylov} to nail to machine $\varepsilon$.
  \item \textbf{$N_Q$ auto-selection}: try several GL orders, keep the
        FP with smallest $\|F\|_\infty$. Different $(\tau, \gamma)$ have
        different sweet spots.
\end{itemize}

\section{Convergence and per-$\gamma$ floor}

\begin{center}
\includegraphics[width=0.95\linewidth]{02_convergence.png}
\end{center}

\begin{center}
\includegraphics[width=0.85\linewidth]{03_floors_per_gamma.png}
\end{center}

\section{Equilibrium statistics vs $\gamma$}

\begin{center}
\includegraphics[width=0.95\linewidth]{01_slope_deficit.png}
\end{center}

\paragraph{Interpretation.}
\begin{itemize}
  \item $\alpha^*$ measures the slope of $\mathrm{logit}\,P$ on the
        aggregate signal $T = \tau \sum_k u_k$. A larger slope means the
        price translates more aggressively into posterior odds: more
        revealing.
  \item $1-R^2$ is the \emph{information deficit}: the fraction of
        $\mathrm{logit}\,P$ variance not explained by $T$. Zero means
        fully revealing (price exactly summarizes $T$); larger means
        partial revelation.
\end{itemize}

As $\gamma$ grows from 0.1 to 5.0, agents become more risk-averse.
The slope $\alpha^*$ and deficit $1-R^2$ trace out the partially-revealing
branch of the K=3 REE.

\section{Price contours $\{P(u_1, u_2, u_3) = p\}$ at each $\gamma$}

Each page shows three slices of the price function $P$ at low / mid / high
$u_3$. The labeled contours are the level sets at $p \in \{0.1, 0.2, ..., 0.9\}$.

\input{contour_pages.tex}

\section{Posterior $\mu(p, u_k)$ tables}

For each $\gamma$, $\mu(p, u_k)$ is the conditional probability $\Pr(v=1 \mid p, u_k)$
that an agent with private signal $u_k$ assigns to the asset paying off, given the
market price $p$. Left panel: heatmap on $(p, u_k)$. Right panel: slices at each $u_k$.

\input{mu_pages.tex}

\section{Logit-$P$ vs $T$ regression}

For each $\gamma$, $\mathrm{logit}\,P$ scatter against $T$, with the
best-fit slope $\alpha^*$ overlaid. Linear fit deviations are the
information deficit.

\begin{center}
\includegraphics[width=0.95\linewidth]{06_logit_scatter.png}
\end{center}

\section{$P$ slice visualizations}

3D price function projected to 2D slices at fixed $u_3$.

\input{slice_pages.tex}

\section{Summary table}

\input{summary_table.tex}

\end{document}
